% --- Archon physics formalization source begin ---
% archon:physics
% archon:covers IPhO2026Problems/problem_ipho_2026_t1_b1.lean
% archon:source-report reports/ipho_2026/problem_ipho_2026_t1_b1.source.json
% archon:problem-id ipho_2026_t1
% archon:part-id T1-B1

\chapter{Ipho Physics problem ipho\_2026\_t1\_b1}
\label{ch:IPhO2026Problems_problem_ipho_2026_t1_b1}

\paragraph{Problem source.}
\#\# Physical scenario

At one instant a positron and an electron, each of mass m and charges of equal
magnitude and opposite sign, are separated by 100*a\_0.  Their velocities are
antiparallel and perpendicular to their separation.  Each particle has angular
momentum of magnitude mu*hbar about the center of mass.  The system is isolated,
classical, non-relativistic, and has only electrostatic interaction.  The Bohr
radius is a\_0 = 4*pi*epsilon\_0*hbar\textasciicircum{}2/(m*e\textasciicircum{}2), and k = 1/(4*pi*epsilon\_0).

\#\# Current subquestion T1-B1

For mu = 4 the pair is bound. Find the maximum electron-positron separation in units of a\_0.

\paragraph{Shared problem context.}
At one instant a positron and an electron, each of mass m and charges of equal
magnitude and opposite sign, are separated by 100*a\_0.  Their velocities are
antiparallel and perpendicular to their separation.  Each particle has angular
momentum of magnitude mu*hbar about the center of mass.  The system is isolated,
classical, non-relativistic, and has only electrostatic interaction.  The Bohr
radius is a\_0 = 4*pi*epsilon\_0*hbar\textasciicircum{}2/(m*e\textasciicircum{}2), and k = 1/(4*pi*epsilon\_0).

\paragraph{Current subquestion.}
For mu = 4 the pair is bound. Find the maximum electron-positron separation in units of a\_0.

\paragraph{Figure/image path.}
ipho\_2026\_source/image/T1\_page-2.png

\paragraph{Formalization target.}
create a compiling Lean file with sorry bodies at `IPhO2026Problems/problem\_ipho\_2026\_t1\_b1.lean`.
The Lean declarations must preserve the physical quantities, dimensions or dimensional roles, figure labels, governing-law hypotheses, and final relation expressed by this problem.
Use Mathlib/Physlib names found through LeanExplore where available. If a domain API is missing, introduce faithful local abstractions rather than scalar placeholder aliases.
This is an answer-blind solve-phase task. Derive a candidate result only from the problem-side evidence above; do not assume a target value, select a tolerance around a desired value, or encode a desired result into a candidate domain.
For numerical questions, define the raw end-to-end quantity and apply an explicit source-derived rounding rule. For identification questions, characterize or prove uniqueness before recording the derived candidate.

\begin{theorem}[Ipho Physics formalization target]
\label{thm:physics:ipho_2026_t1_b1:target}
\lean{IPhO2026T1B1.max_separation_eq, IPhO2026T1B1.max_separation_over_bohrRadius}
\leanok
For $\mu = 4$ the pair is bound, and the maximum electron--positron
separation along the Coulomb orbit is
$r_{\max} = p/(1-\varepsilon) = 128\,a_0/(1 - 7/25) = (1600/9)\,a_0$,
i.e. $r_{\max}/a_0 = 1600/9 \approx 177.78$.
\end{theorem}
\begin{proof}
\leanok
Kernel-checked in Lean as \texttt{IPhO2026T1B1.max\_separation\_eq} and
\texttt{IPhO2026T1B1.max\_separation\_over\_bohrRadius} (no \texttt{sorry}).
Derivation chain from problem-side data only: $L = 2\mu\hbar = 8\hbar$;
initial energy $E = L^2/(m r_0^2) - k e^2/r_0
= -(36/10000)\,\hbar^2/(m a_0^2) < 0$ (the pair is bound); Hint~1 gives
$\varepsilon = \sqrt{1 + 4L^2E/(k^2e^4m)} = \sqrt{49/625} = 7/25$; the
Kepler/Coulomb semi-latus-rectum relation gives
$p = L^2/((m/2)\,k e^2) = 128\,a_0$; the apoapsis of the conic
$r(\theta) = p/(1 - \varepsilon\cos\theta)$ (Hint~2, supremum over $\theta$,
attained at $\cos\theta = 1$) is $p/(1-\varepsilon) = (1600/9)\,a_0$.
Consistency with Fig.~1b: the periapsis is $p/(1+\varepsilon) = 100\,a_0 = r_0$.
\end{proof}
% --- Archon physics formalization source end ---
